% Options for packages loaded elsewhere
\PassOptionsToPackage{unicode}{hyperref}
\PassOptionsToPackage{hyphens}{url}
\PassOptionsToPackage{dvipsnames,svgnames,x11names}{xcolor}
%
\documentclass[
  letterpaper,
  DIV=11,
  numbers=noendperiod]{scrartcl}

\usepackage{amsmath,amssymb}
\usepackage{iftex}
\ifPDFTeX
  \usepackage[T1]{fontenc}
  \usepackage[utf8]{inputenc}
  \usepackage{textcomp} % provide euro and other symbols
\else % if luatex or xetex
  \usepackage{unicode-math}
  \defaultfontfeatures{Scale=MatchLowercase}
  \defaultfontfeatures[\rmfamily]{Ligatures=TeX,Scale=1}
\fi
\usepackage{lmodern}
\ifPDFTeX\else  
    % xetex/luatex font selection
\fi
% Use upquote if available, for straight quotes in verbatim environments
\IfFileExists{upquote.sty}{\usepackage{upquote}}{}
\IfFileExists{microtype.sty}{% use microtype if available
  \usepackage[]{microtype}
  \UseMicrotypeSet[protrusion]{basicmath} % disable protrusion for tt fonts
}{}
\makeatletter
\@ifundefined{KOMAClassName}{% if non-KOMA class
  \IfFileExists{parskip.sty}{%
    \usepackage{parskip}
  }{% else
    \setlength{\parindent}{0pt}
    \setlength{\parskip}{6pt plus 2pt minus 1pt}}
}{% if KOMA class
  \KOMAoptions{parskip=half}}
\makeatother
\usepackage{xcolor}
\setlength{\emergencystretch}{3em} % prevent overfull lines
\setcounter{secnumdepth}{-\maxdimen} % remove section numbering
% Make \paragraph and \subparagraph free-standing
\makeatletter
\ifx\paragraph\undefined\else
  \let\oldparagraph\paragraph
  \renewcommand{\paragraph}{
    \@ifstar
      \xxxParagraphStar
      \xxxParagraphNoStar
  }
  \newcommand{\xxxParagraphStar}[1]{\oldparagraph*{#1}\mbox{}}
  \newcommand{\xxxParagraphNoStar}[1]{\oldparagraph{#1}\mbox{}}
\fi
\ifx\subparagraph\undefined\else
  \let\oldsubparagraph\subparagraph
  \renewcommand{\subparagraph}{
    \@ifstar
      \xxxSubParagraphStar
      \xxxSubParagraphNoStar
  }
  \newcommand{\xxxSubParagraphStar}[1]{\oldsubparagraph*{#1}\mbox{}}
  \newcommand{\xxxSubParagraphNoStar}[1]{\oldsubparagraph{#1}\mbox{}}
\fi
\makeatother

\usepackage{color}
\usepackage{fancyvrb}
\newcommand{\VerbBar}{|}
\newcommand{\VERB}{\Verb[commandchars=\\\{\}]}
\DefineVerbatimEnvironment{Highlighting}{Verbatim}{commandchars=\\\{\}}
% Add ',fontsize=\small' for more characters per line
\usepackage{framed}
\definecolor{shadecolor}{RGB}{241,243,245}
\newenvironment{Shaded}{\begin{snugshade}}{\end{snugshade}}
\newcommand{\AlertTok}[1]{\textcolor[rgb]{0.68,0.00,0.00}{#1}}
\newcommand{\AnnotationTok}[1]{\textcolor[rgb]{0.37,0.37,0.37}{#1}}
\newcommand{\AttributeTok}[1]{\textcolor[rgb]{0.40,0.45,0.13}{#1}}
\newcommand{\BaseNTok}[1]{\textcolor[rgb]{0.68,0.00,0.00}{#1}}
\newcommand{\BuiltInTok}[1]{\textcolor[rgb]{0.00,0.23,0.31}{#1}}
\newcommand{\CharTok}[1]{\textcolor[rgb]{0.13,0.47,0.30}{#1}}
\newcommand{\CommentTok}[1]{\textcolor[rgb]{0.37,0.37,0.37}{#1}}
\newcommand{\CommentVarTok}[1]{\textcolor[rgb]{0.37,0.37,0.37}{\textit{#1}}}
\newcommand{\ConstantTok}[1]{\textcolor[rgb]{0.56,0.35,0.01}{#1}}
\newcommand{\ControlFlowTok}[1]{\textcolor[rgb]{0.00,0.23,0.31}{\textbf{#1}}}
\newcommand{\DataTypeTok}[1]{\textcolor[rgb]{0.68,0.00,0.00}{#1}}
\newcommand{\DecValTok}[1]{\textcolor[rgb]{0.68,0.00,0.00}{#1}}
\newcommand{\DocumentationTok}[1]{\textcolor[rgb]{0.37,0.37,0.37}{\textit{#1}}}
\newcommand{\ErrorTok}[1]{\textcolor[rgb]{0.68,0.00,0.00}{#1}}
\newcommand{\ExtensionTok}[1]{\textcolor[rgb]{0.00,0.23,0.31}{#1}}
\newcommand{\FloatTok}[1]{\textcolor[rgb]{0.68,0.00,0.00}{#1}}
\newcommand{\FunctionTok}[1]{\textcolor[rgb]{0.28,0.35,0.67}{#1}}
\newcommand{\ImportTok}[1]{\textcolor[rgb]{0.00,0.46,0.62}{#1}}
\newcommand{\InformationTok}[1]{\textcolor[rgb]{0.37,0.37,0.37}{#1}}
\newcommand{\KeywordTok}[1]{\textcolor[rgb]{0.00,0.23,0.31}{\textbf{#1}}}
\newcommand{\NormalTok}[1]{\textcolor[rgb]{0.00,0.23,0.31}{#1}}
\newcommand{\OperatorTok}[1]{\textcolor[rgb]{0.37,0.37,0.37}{#1}}
\newcommand{\OtherTok}[1]{\textcolor[rgb]{0.00,0.23,0.31}{#1}}
\newcommand{\PreprocessorTok}[1]{\textcolor[rgb]{0.68,0.00,0.00}{#1}}
\newcommand{\RegionMarkerTok}[1]{\textcolor[rgb]{0.00,0.23,0.31}{#1}}
\newcommand{\SpecialCharTok}[1]{\textcolor[rgb]{0.37,0.37,0.37}{#1}}
\newcommand{\SpecialStringTok}[1]{\textcolor[rgb]{0.13,0.47,0.30}{#1}}
\newcommand{\StringTok}[1]{\textcolor[rgb]{0.13,0.47,0.30}{#1}}
\newcommand{\VariableTok}[1]{\textcolor[rgb]{0.07,0.07,0.07}{#1}}
\newcommand{\VerbatimStringTok}[1]{\textcolor[rgb]{0.13,0.47,0.30}{#1}}
\newcommand{\WarningTok}[1]{\textcolor[rgb]{0.37,0.37,0.37}{\textit{#1}}}

\providecommand{\tightlist}{%
  \setlength{\itemsep}{0pt}\setlength{\parskip}{0pt}}\usepackage{longtable,booktabs,array}
\usepackage{calc} % for calculating minipage widths
% Correct order of tables after \paragraph or \subparagraph
\usepackage{etoolbox}
\makeatletter
\patchcmd\longtable{\par}{\if@noskipsec\mbox{}\fi\par}{}{}
\makeatother
% Allow footnotes in longtable head/foot
\IfFileExists{footnotehyper.sty}{\usepackage{footnotehyper}}{\usepackage{footnote}}
\makesavenoteenv{longtable}
\usepackage{graphicx}
\makeatletter
\newsavebox\pandoc@box
\newcommand*\pandocbounded[1]{% scales image to fit in text height/width
  \sbox\pandoc@box{#1}%
  \Gscale@div\@tempa{\textheight}{\dimexpr\ht\pandoc@box+\dp\pandoc@box\relax}%
  \Gscale@div\@tempb{\linewidth}{\wd\pandoc@box}%
  \ifdim\@tempb\p@<\@tempa\p@\let\@tempa\@tempb\fi% select the smaller of both
  \ifdim\@tempa\p@<\p@\scalebox{\@tempa}{\usebox\pandoc@box}%
  \else\usebox{\pandoc@box}%
  \fi%
}
% Set default figure placement to htbp
\def\fps@figure{htbp}
\makeatother

\KOMAoption{captions}{tableheading}
\makeatletter
\@ifpackageloaded{caption}{}{\usepackage{caption}}
\AtBeginDocument{%
\ifdefined\contentsname
  \renewcommand*\contentsname{Table of contents}
\else
  \newcommand\contentsname{Table of contents}
\fi
\ifdefined\listfigurename
  \renewcommand*\listfigurename{List of Figures}
\else
  \newcommand\listfigurename{List of Figures}
\fi
\ifdefined\listtablename
  \renewcommand*\listtablename{List of Tables}
\else
  \newcommand\listtablename{List of Tables}
\fi
\ifdefined\figurename
  \renewcommand*\figurename{Figure}
\else
  \newcommand\figurename{Figure}
\fi
\ifdefined\tablename
  \renewcommand*\tablename{Table}
\else
  \newcommand\tablename{Table}
\fi
}
\@ifpackageloaded{float}{}{\usepackage{float}}
\floatstyle{ruled}
\@ifundefined{c@chapter}{\newfloat{codelisting}{h}{lop}}{\newfloat{codelisting}{h}{lop}[chapter]}
\floatname{codelisting}{Listing}
\newcommand*\listoflistings{\listof{codelisting}{List of Listings}}
\makeatother
\makeatletter
\makeatother
\makeatletter
\@ifpackageloaded{caption}{}{\usepackage{caption}}
\@ifpackageloaded{subcaption}{}{\usepackage{subcaption}}
\makeatother

\usepackage{bookmark}

\IfFileExists{xurl.sty}{\usepackage{xurl}}{} % add URL line breaks if available
\urlstyle{same} % disable monospaced font for URLs
\hypersetup{
  pdftitle={FLUME Patent Figures},
  colorlinks=true,
  linkcolor={blue},
  filecolor={Maroon},
  citecolor={Blue},
  urlcolor={Blue},
  pdfcreator={LaTeX via pandoc}}


\title{FLUME Patent Figures}
\author{}
\date{}

\begin{document}
\maketitle


\section{FLUME Patent Figures}\label{flume-patent-figures}

\subsection{FIG. 1: System Architecture - Triune Hierarchical
Compression
Pipeline}\label{fig.-1-system-architecture---triune-hierarchical-compression-pipeline}

\begin{Shaded}
\begin{Highlighting}[]
\NormalTok{graph LR}
\NormalTok{    subgraph Knower["Knower Encoder (2048D)"]}
\NormalTok{        A[Text Input] {-}{-}\textgreater{} B[LLM Embedding]}
\NormalTok{        B {-}{-}\textgreater{} C[2048D Vector]}
\NormalTok{    end}
    
\NormalTok{    subgraph Thinker["Thinker VAE (512D)"]}
\NormalTok{        C {-}{-}\textgreater{} D[Encoder]}
\NormalTok{        D {-}{-}\textgreater{} E[μ, log σ²]}
\NormalTok{        E {-}{-}\textgreater{} F[Reparameterization]}
\NormalTok{        F {-}{-}\textgreater{} G[512D Latent z]}
\NormalTok{    end}
    
\NormalTok{    subgraph Doer["Doer Projector (12D)"]}
\NormalTok{        G {-}{-}\textgreater{} H[Projection Layer]}
\NormalTok{        H {-}{-}\textgreater{} I[12D Observable]}
\NormalTok{    end}
    
\NormalTok{    subgraph Coherence["Coherence Regularizer"]}
\NormalTok{        G {-}{-}\textgreater{} J[Coherence Loss]}
\NormalTok{        J {-}{-}\textgreater{} K[0.5 Target]}
\NormalTok{    end}
    
\NormalTok{    subgraph Navigator["Trajectory Navigator"]}
\NormalTok{        G {-}{-}\textgreater{} L[Geodesic Prediction]}
\NormalTok{        L {-}{-}\textgreater{} M[Trajectory Points]}
\NormalTok{    end}
    
\NormalTok{    I {-}{-}\textgreater{} N[Observable Output]}
\NormalTok{    M {-}{-}\textgreater{} O[Journey Logging]}
    
\NormalTok{    style Knower fill:lightblue,stroke:blue,stroke{-}width:2px}
\NormalTok{    style Thinker fill:lightgreen,stroke:green,stroke{-}width:2px}
\NormalTok{    style Doer fill:lightcoral,stroke:red,stroke{-}width:2px}
\NormalTok{    style Coherence fill:yellow,stroke:orange,stroke{-}width:2px,stroke{-}dasharray: 5 5}
\NormalTok{    style Navigator fill:lavender,stroke:purple,stroke{-}width:2px,stroke{-}dasharray: 5 5}
\end{Highlighting}
\end{Shaded}

\textbf{Description:} System architecture showing Knower Encoder (2048D)
→ Thinker VAE (512D) → Doer Projector (12D) flow with Coherence
Regularizer (0.5 target) and Trajectory Navigator.

\begin{center}\rule{0.5\linewidth}{0.5pt}\end{center}

\subsection{FIG. 2: VAE Encoder-Decoder with HIHO Coherence
Loss}\label{fig.-2-vae-encoder-decoder-with-hiho-coherence-loss}

\begin{figure}[H]

{\centering \pandocbounded{\includegraphics[keepaspectratio]{figures_files/mediabag/svg/fig02_vae.pdf}}

}

\caption{FIG. 2}

\end{figure}%

\textbf{Description:} VAE architecture showing encoder producing μ and
log(σ²), reparameterization trick (z = μ + ε·σ), decoder reconstruction,
and HIHO coherence loss toward 0.5 target.

\begin{center}\rule{0.5\linewidth}{0.5pt}\end{center}

\subsection{FIG. 3: 12D Physics-Grounded State (Smith's 4
Fabrics)}\label{fig.-3-12d-physics-grounded-state-smiths-4-fabrics}

\begin{figure}[H]

{\centering \pandocbounded{\includegraphics[keepaspectratio]{figures_files/mediabag/svg/fig03_12d_state.pdf}}

}

\caption{FIG. 3}

\end{figure}%

\textbf{Description:} 12D observable state organized as four fabrics of
three dimensions each: Space (x,y,z), Field (Tempic,Electric,Magnetic),
Control (Rotation,Precession,Charge), Precipitation
(Awareness,Novelty,Manifestation).

\begin{center}\rule{0.5\linewidth}{0.5pt}\end{center}

\subsection{FIG. 4: Continuous Trajectory Prediction in 512D
Manifold}\label{fig.-4-continuous-trajectory-prediction-in-512d-manifold}

\begin{figure}[H]

{\centering \pandocbounded{\includegraphics[keepaspectratio]{figures_files/mediabag/svg/fig04_trajectory.pdf}}

}

\caption{FIG. 4}

\end{figure}%

\textbf{Description:} Geodesic trajectory through 512D manifold from
start latent to goal latent, with interpolation (z = α·z₁ + (1-α)·z₂)
and projection to 12D observable state.

\begin{center}\rule{0.5\linewidth}{0.5pt}\end{center}

\subsection{FIG. 5: HIHO Double-Well Potential Energy
Landscape}\label{fig.-5-hiho-double-well-potential-energy-landscape}

\begin{Shaded}
\begin{Highlighting}[]
\CommentTok{\#| label: fig{-}hiho}
\CommentTok{\#| fig{-}cap: "HIHO Double{-}Well Potential"}
\ImportTok{import}\NormalTok{ matplotlib.pyplot }\ImportTok{as}\NormalTok{ plt}
\ImportTok{import}\NormalTok{ numpy }\ImportTok{as}\NormalTok{ np}

\CommentTok{\# Set publication quality}
\NormalTok{plt.rcParams[}\StringTok{\textquotesingle{}figure.dpi\textquotesingle{}}\NormalTok{] }\OperatorTok{=} \DecValTok{300}
\NormalTok{plt.rcParams[}\StringTok{\textquotesingle{}font.size\textquotesingle{}}\NormalTok{] }\OperatorTok{=} \DecValTok{10}

\CommentTok{\# Double{-}well potential: V(x) = (x {-} 0.5)\^{}4 {-} 0.5*(x {-} 0.5)\^{}2}
\NormalTok{x }\OperatorTok{=}\NormalTok{ np.linspace(}\DecValTok{0}\NormalTok{, }\DecValTok{1}\NormalTok{, }\DecValTok{100}\NormalTok{)}
\NormalTok{V }\OperatorTok{=}\NormalTok{ (x }\OperatorTok{{-}} \FloatTok{0.5}\NormalTok{)}\OperatorTok{**}\DecValTok{4} \OperatorTok{{-}} \FloatTok{0.5} \OperatorTok{*}\NormalTok{ (x }\OperatorTok{{-}} \FloatTok{0.5}\NormalTok{)}\OperatorTok{**}\DecValTok{2}
\NormalTok{V }\OperatorTok{=}\NormalTok{ V }\OperatorTok{{-}}\NormalTok{ V.}\BuiltInTok{min}\NormalTok{()  }\CommentTok{\# Shift to make minimum at 0}
\NormalTok{V }\OperatorTok{=}\NormalTok{ V }\OperatorTok{/}\NormalTok{ V.}\BuiltInTok{max}\NormalTok{() }\OperatorTok{*} \DecValTok{2}  \CommentTok{\# Scale}

\NormalTok{fig, ax }\OperatorTok{=}\NormalTok{ plt.subplots(}\DecValTok{1}\NormalTok{, }\DecValTok{1}\NormalTok{, figsize}\OperatorTok{=}\NormalTok{(}\DecValTok{8}\NormalTok{, }\DecValTok{5}\NormalTok{))}
\NormalTok{ax.plot(x, V, }\StringTok{\textquotesingle{}b{-}\textquotesingle{}}\NormalTok{, linewidth}\OperatorTok{=}\DecValTok{3}\NormalTok{, label}\OperatorTok{=}\StringTok{\textquotesingle{}Free Energy Landscape\textquotesingle{}}\NormalTok{)}

\CommentTok{\# Minimum at 0.5}
\NormalTok{ax.axvline(x}\OperatorTok{=}\FloatTok{0.5}\NormalTok{, color}\OperatorTok{=}\StringTok{\textquotesingle{}red\textquotesingle{}}\NormalTok{, linestyle}\OperatorTok{=}\StringTok{\textquotesingle{}{-}{-}\textquotesingle{}}\NormalTok{, linewidth}\OperatorTok{=}\DecValTok{2}\NormalTok{, label}\OperatorTok{=}\StringTok{\textquotesingle{}HIHO Target (0.5)\textquotesingle{}}\NormalTok{)}
\NormalTok{ax.plot(}\FloatTok{0.5}\NormalTok{, }\DecValTok{0}\NormalTok{, }\StringTok{\textquotesingle{}ro\textquotesingle{}}\NormalTok{, markersize}\OperatorTok{=}\DecValTok{15}\NormalTok{)}
\NormalTok{ax.text(}\FloatTok{0.5}\NormalTok{, }\FloatTok{0.2}\NormalTok{, }\StringTok{"Minimum}\CharTok{\textbackslash{}n}\StringTok{(0.5)"}\NormalTok{, ha}\OperatorTok{=}\StringTok{\textquotesingle{}center\textquotesingle{}}\NormalTok{, fontsize}\OperatorTok{=}\DecValTok{9}\NormalTok{, fontweight}\OperatorTok{=}\StringTok{\textquotesingle{}bold\textquotesingle{}}\NormalTok{)}

\CommentTok{\# Wells}
\NormalTok{ax.fill\_between(x, }\DecValTok{0}\NormalTok{, V, alpha}\OperatorTok{=}\FloatTok{0.3}\NormalTok{, color}\OperatorTok{=}\StringTok{\textquotesingle{}blue\textquotesingle{}}\NormalTok{)}

\CommentTok{\# Annotations}
\NormalTok{ax.text(}\FloatTok{0.25}\NormalTok{, }\FloatTok{1.5}\NormalTok{, }\StringTok{"Exploration}\CharTok{\textbackslash{}n}\StringTok{(novelty)"}\NormalTok{, ha}\OperatorTok{=}\StringTok{\textquotesingle{}center\textquotesingle{}}\NormalTok{, fontsize}\OperatorTok{=}\DecValTok{9}\NormalTok{, }
\NormalTok{        bbox}\OperatorTok{=}\BuiltInTok{dict}\NormalTok{(boxstyle}\OperatorTok{=}\StringTok{\textquotesingle{}round\textquotesingle{}}\NormalTok{, facecolor}\OperatorTok{=}\StringTok{\textquotesingle{}lightblue\textquotesingle{}}\NormalTok{, alpha}\OperatorTok{=}\FloatTok{0.5}\NormalTok{))}
\NormalTok{ax.text(}\FloatTok{0.75}\NormalTok{, }\FloatTok{1.5}\NormalTok{, }\StringTok{"Exploitation}\CharTok{\textbackslash{}n}\StringTok{(precipitation)"}\NormalTok{, ha}\OperatorTok{=}\StringTok{\textquotesingle{}center\textquotesingle{}}\NormalTok{, fontsize}\OperatorTok{=}\DecValTok{9}\NormalTok{,}
\NormalTok{        bbox}\OperatorTok{=}\BuiltInTok{dict}\NormalTok{(boxstyle}\OperatorTok{=}\StringTok{\textquotesingle{}round\textquotesingle{}}\NormalTok{, facecolor}\OperatorTok{=}\StringTok{\textquotesingle{}lightcoral\textquotesingle{}}\NormalTok{, alpha}\OperatorTok{=}\FloatTok{0.5}\NormalTok{))}

\CommentTok{\# Thermodynamic derivation}
\NormalTok{ax.text(}\FloatTok{0.5}\NormalTok{, }\FloatTok{2.5}\NormalTok{, }\StringTok{"Thermodynamic Ground State:}\CharTok{\textbackslash{}n}\StringTok{Max Entropy at p = 0.5"}\NormalTok{, }
\NormalTok{        ha}\OperatorTok{=}\StringTok{\textquotesingle{}center\textquotesingle{}}\NormalTok{, fontsize}\OperatorTok{=}\DecValTok{9}\NormalTok{, style}\OperatorTok{=}\StringTok{\textquotesingle{}italic\textquotesingle{}}\NormalTok{,}
\NormalTok{        bbox}\OperatorTok{=}\BuiltInTok{dict}\NormalTok{(boxstyle}\OperatorTok{=}\StringTok{\textquotesingle{}round\textquotesingle{}}\NormalTok{, facecolor}\OperatorTok{=}\StringTok{\textquotesingle{}yellow\textquotesingle{}}\NormalTok{, alpha}\OperatorTok{=}\FloatTok{0.5}\NormalTok{))}

\NormalTok{ax.set\_xlabel(}\StringTok{"Coherence"}\NormalTok{, fontsize}\OperatorTok{=}\DecValTok{11}\NormalTok{)}
\NormalTok{ax.set\_ylabel(}\StringTok{"Free Energy"}\NormalTok{, fontsize}\OperatorTok{=}\DecValTok{11}\NormalTok{)}
\NormalTok{ax.set\_title(}\StringTok{"FIG. 5: HIHO Double{-}Well Potential"}\NormalTok{, fontsize}\OperatorTok{=}\DecValTok{12}\NormalTok{, fontweight}\OperatorTok{=}\StringTok{\textquotesingle{}bold\textquotesingle{}}\NormalTok{)}
\NormalTok{ax.legend(loc}\OperatorTok{=}\StringTok{\textquotesingle{}upper right\textquotesingle{}}\NormalTok{, fontsize}\OperatorTok{=}\DecValTok{9}\NormalTok{)}
\NormalTok{ax.grid(}\VariableTok{True}\NormalTok{, alpha}\OperatorTok{=}\FloatTok{0.3}\NormalTok{)}
\NormalTok{plt.tight\_layout()}
\NormalTok{plt.show()}
\end{Highlighting}
\end{Shaded}

\textbf{Description:} Free energy landscape showing double-well
potential with minimum at 0.5 coherence (thermodynamic ground state,
maximum entropy).

\begin{center}\rule{0.5\linewidth}{0.5pt}\end{center}

\subsection{FIG. 6: Training Loss Convergence with Coherence
Regularization}\label{fig.-6-training-loss-convergence-with-coherence-regularization}

\begin{Shaded}
\begin{Highlighting}[]
\CommentTok{\#| label: fig{-}training}
\CommentTok{\#| fig{-}cap: "Training Loss Convergence"}
\ImportTok{import}\NormalTok{ matplotlib.pyplot }\ImportTok{as}\NormalTok{ plt}
\ImportTok{import}\NormalTok{ numpy }\ImportTok{as}\NormalTok{ np}

\CommentTok{\# Simulated training curves}
\NormalTok{epochs }\OperatorTok{=}\NormalTok{ np.arange(}\DecValTok{1}\NormalTok{, }\DecValTok{51}\NormalTok{)}
\NormalTok{mse\_loss }\OperatorTok{=} \FloatTok{0.8} \OperatorTok{*}\NormalTok{ np.exp(}\OperatorTok{{-}}\FloatTok{0.08} \OperatorTok{*}\NormalTok{ epochs) }\OperatorTok{+} \FloatTok{0.13}
\NormalTok{kl\_loss }\OperatorTok{=} \FloatTok{0.6} \OperatorTok{*}\NormalTok{ np.exp(}\OperatorTok{{-}}\FloatTok{0.06} \OperatorTok{*}\NormalTok{ epochs) }\OperatorTok{+} \FloatTok{0.43}
\NormalTok{coherence\_loss }\OperatorTok{=} \FloatTok{0.4} \OperatorTok{*}\NormalTok{ np.exp(}\OperatorTok{{-}}\FloatTok{0.1} \OperatorTok{*}\NormalTok{ epochs) }\OperatorTok{+} \FloatTok{0.15}
\NormalTok{total\_loss }\OperatorTok{=}\NormalTok{ mse\_loss }\OperatorTok{+} \FloatTok{0.1} \OperatorTok{*}\NormalTok{ kl\_loss }\OperatorTok{+} \FloatTok{0.1} \OperatorTok{*}\NormalTok{ coherence\_loss}

\NormalTok{fig, ax }\OperatorTok{=}\NormalTok{ plt.subplots(}\DecValTok{1}\NormalTok{, }\DecValTok{1}\NormalTok{, figsize}\OperatorTok{=}\NormalTok{(}\DecValTok{8}\NormalTok{, }\DecValTok{5}\NormalTok{))}
\NormalTok{ax.plot(epochs, mse\_loss, }\StringTok{\textquotesingle{}b{-}\textquotesingle{}}\NormalTok{, linewidth}\OperatorTok{=}\DecValTok{2}\NormalTok{, label}\OperatorTok{=}\StringTok{\textquotesingle{}Reconstruction (MSE)\textquotesingle{}}\NormalTok{)}
\NormalTok{ax.plot(epochs, kl\_loss, }\StringTok{\textquotesingle{}g{-}\textquotesingle{}}\NormalTok{, linewidth}\OperatorTok{=}\DecValTok{2}\NormalTok{, label}\OperatorTok{=}\StringTok{\textquotesingle{}KL Divergence\textquotesingle{}}\NormalTok{)}
\NormalTok{ax.plot(epochs, coherence\_loss, }\StringTok{\textquotesingle{}r{-}\textquotesingle{}}\NormalTok{, linewidth}\OperatorTok{=}\DecValTok{2}\NormalTok{, label}\OperatorTok{=}\StringTok{\textquotesingle{}Coherence Loss\textquotesingle{}}\NormalTok{)}
\NormalTok{ax.plot(epochs, total\_loss, }\StringTok{\textquotesingle{}k{-}\textquotesingle{}}\NormalTok{, linewidth}\OperatorTok{=}\DecValTok{3}\NormalTok{, label}\OperatorTok{=}\StringTok{\textquotesingle{}Total Loss\textquotesingle{}}\NormalTok{)}

\CommentTok{\# Annotations}
\NormalTok{ax.axhline(y}\OperatorTok{=}\FloatTok{0.1322}\NormalTok{, color}\OperatorTok{=}\StringTok{\textquotesingle{}blue\textquotesingle{}}\NormalTok{, linestyle}\OperatorTok{=}\StringTok{\textquotesingle{}{-}{-}\textquotesingle{}}\NormalTok{, linewidth}\OperatorTok{=}\DecValTok{1}\NormalTok{, alpha}\OperatorTok{=}\FloatTok{0.5}\NormalTok{)}
\NormalTok{ax.text(}\DecValTok{45}\NormalTok{, }\FloatTok{0.14}\NormalTok{, }\StringTok{"0.1322"}\NormalTok{, ha}\OperatorTok{=}\StringTok{\textquotesingle{}right\textquotesingle{}}\NormalTok{, fontsize}\OperatorTok{=}\DecValTok{9}\NormalTok{, color}\OperatorTok{=}\StringTok{\textquotesingle{}blue\textquotesingle{}}\NormalTok{)}

\NormalTok{ax.axhline(y}\OperatorTok{=}\FloatTok{0.4329}\NormalTok{, color}\OperatorTok{=}\StringTok{\textquotesingle{}green\textquotesingle{}}\NormalTok{, linestyle}\OperatorTok{=}\StringTok{\textquotesingle{}{-}{-}\textquotesingle{}}\NormalTok{, linewidth}\OperatorTok{=}\DecValTok{1}\NormalTok{, alpha}\OperatorTok{=}\FloatTok{0.5}\NormalTok{)}
\NormalTok{ax.text(}\DecValTok{45}\NormalTok{, }\FloatTok{0.44}\NormalTok{, }\StringTok{"0.4329"}\NormalTok{, ha}\OperatorTok{=}\StringTok{\textquotesingle{}right\textquotesingle{}}\NormalTok{, fontsize}\OperatorTok{=}\DecValTok{9}\NormalTok{, color}\OperatorTok{=}\StringTok{\textquotesingle{}green\textquotesingle{}}\NormalTok{)}

\NormalTok{ax.axhline(y}\OperatorTok{=}\FloatTok{0.63}\NormalTok{, color}\OperatorTok{=}\StringTok{\textquotesingle{}red\textquotesingle{}}\NormalTok{, linestyle}\OperatorTok{=}\StringTok{\textquotesingle{}{-}{-}\textquotesingle{}}\NormalTok{, linewidth}\OperatorTok{=}\DecValTok{1}\NormalTok{, alpha}\OperatorTok{=}\FloatTok{0.5}\NormalTok{)}
\NormalTok{ax.text(}\DecValTok{45}\NormalTok{, }\FloatTok{0.64}\NormalTok{, }\StringTok{"0.63 (mean coherence)"}\NormalTok{, ha}\OperatorTok{=}\StringTok{\textquotesingle{}right\textquotesingle{}}\NormalTok{, fontsize}\OperatorTok{=}\DecValTok{9}\NormalTok{, color}\OperatorTok{=}\StringTok{\textquotesingle{}red\textquotesingle{}}\NormalTok{)}

\CommentTok{\# Final values box}
\NormalTok{box }\OperatorTok{=}\NormalTok{ plt.Rectangle((}\DecValTok{25}\NormalTok{, }\FloatTok{0.8}\NormalTok{), }\DecValTok{20}\NormalTok{, }\FloatTok{0.35}\NormalTok{, boxstyle}\OperatorTok{=}\StringTok{"round,pad=0.1"}\NormalTok{,}
\NormalTok{                    facecolor}\OperatorTok{=}\StringTok{\textquotesingle{}yellow\textquotesingle{}}\NormalTok{, edgecolor}\OperatorTok{=}\StringTok{\textquotesingle{}orange\textquotesingle{}}\NormalTok{, linewidth}\OperatorTok{=}\DecValTok{2}\NormalTok{, alpha}\OperatorTok{=}\FloatTok{0.5}\NormalTok{)}
\NormalTok{ax.add\_patch(box)}
\NormalTok{ax.text(}\DecValTok{35}\NormalTok{, }\FloatTok{1.05}\NormalTok{, }\StringTok{"Final (Epoch 50):}\CharTok{\textbackslash{}n}\StringTok{MSE: 0.1322, KL: 0.4329, Coherence: 0.63"}\NormalTok{,}
\NormalTok{        ha}\OperatorTok{=}\StringTok{\textquotesingle{}center\textquotesingle{}}\NormalTok{, fontsize}\OperatorTok{=}\DecValTok{9}\NormalTok{, fontweight}\OperatorTok{=}\StringTok{\textquotesingle{}bold\textquotesingle{}}\NormalTok{)}

\NormalTok{ax.set\_xlim(}\DecValTok{0}\NormalTok{, }\DecValTok{50}\NormalTok{)}
\NormalTok{ax.set\_ylim(}\DecValTok{0}\NormalTok{, }\FloatTok{1.2}\NormalTok{)}
\NormalTok{ax.set\_xlabel(}\StringTok{"Training Epochs"}\NormalTok{, fontsize}\OperatorTok{=}\DecValTok{11}\NormalTok{)}
\NormalTok{ax.set\_ylabel(}\StringTok{"Loss"}\NormalTok{, fontsize}\OperatorTok{=}\DecValTok{11}\NormalTok{)}
\NormalTok{ax.set\_title(}\StringTok{"FIG. 6: Training Loss Convergence"}\NormalTok{, fontsize}\OperatorTok{=}\DecValTok{12}\NormalTok{, fontweight}\OperatorTok{=}\StringTok{\textquotesingle{}bold\textquotesingle{}}\NormalTok{)}
\NormalTok{ax.legend(loc}\OperatorTok{=}\StringTok{\textquotesingle{}upper right\textquotesingle{}}\NormalTok{, fontsize}\OperatorTok{=}\DecValTok{9}\NormalTok{)}
\NormalTok{ax.grid(}\VariableTok{True}\NormalTok{, alpha}\OperatorTok{=}\FloatTok{0.3}\NormalTok{)}
\NormalTok{plt.tight\_layout()}
\NormalTok{plt.show()}
\end{Highlighting}
\end{Shaded}

\textbf{Description:} Training loss curves showing convergence of
reconstruction (MSE), KL divergence, and coherence loss over 50 epochs,
with final values: MSE=0.1322, KL=0.4329, Coherence=0.63.

\begin{center}\rule{0.5\linewidth}{0.5pt}\end{center}

\subsection{FIG. 7: Journey Tracking Dual-Tier
Logging}\label{fig.-7-journey-tracking-dual-tier-logging}

\begin{Shaded}
\begin{Highlighting}[]
\NormalTok{graph TD}
\NormalTok{    subgraph Observable["12D Observable State Tier"]}
\NormalTok{        O1[spatial\_x]}
\NormalTok{        O2[tempic]}
\NormalTok{        O3[logic]}
\NormalTok{        O4[awareness]}
\NormalTok{        O5[coherence]}
\NormalTok{    end}
    
\NormalTok{    subgraph Semantic["2048D Semantic Context Tier"]}
\NormalTok{        S1[LLM Embedding]}
\NormalTok{        S2[Semantic Intent]}
\NormalTok{        S3[Context Vector]}
\NormalTok{    end}
    
\NormalTok{    subgraph Logging["Per{-}Step Logging"]}
\NormalTok{        L1[Coherence Tracking]}
\NormalTok{        L2[Phi Score]}
\NormalTok{        L3[Thermodynamic State]}
\NormalTok{        L4[Topological Features]}
\NormalTok{    end}
    
\NormalTok{    Semantic {-}{-}\textgreater{}|Projection| Observable}
\NormalTok{    Observable {-}{-}\textgreater{}|Per{-}Step| Logging}
\NormalTok{    Logging {-}{-}\textgreater{} E[Journey Export]}
    
\NormalTok{    style Observable fill:lightcoral,stroke:red,stroke{-}width:2px}
\NormalTok{    style Semantic fill:lightblue,stroke:blue,stroke{-}width:2px}
\NormalTok{    style Logging fill:yellow,stroke:orange,stroke{-}width:2px}
\end{Highlighting}
\end{Shaded}

\textbf{Description:} Dual-tier journey tracking with 12D Observable
State Tier and 2048D Semantic Context Tier, per-step logging, coherence
tracking, phi score computation, thermodynamic state, topological
features, and journey export.

\begin{center}\rule{0.5\linewidth}{0.5pt}\end{center}

\subsection{FIG. 8: Multi-Scale Reasoning
Flowchart}\label{fig.-8-multi-scale-reasoning-flowchart}

\begin{Shaded}
\begin{Highlighting}[]
\NormalTok{graph TD}
\NormalTok{    Start([Start]) {-}{-}\textgreater{} Task[Task: Natural Language]}
    
\NormalTok{    subgraph MultiScale["Multi{-}Scale Operation"]}
\NormalTok{        Knower["Knower Scale (2048D)\textless{}br/\textgreater{}Exhaustive Search"]}
\NormalTok{        Thinker["Thinker Scale (512D)\textless{}br/\textgreater{}Trajectory Prediction"]}
\NormalTok{        Doer["Doer Scale (12D)\textless{}br/\textgreater{}Physical Grounding"]}
\NormalTok{    end}
    
\NormalTok{    Task {-}{-}\textgreater{} Knower}
\NormalTok{    Knower {-}{-}\textgreater{} Thinker}
\NormalTok{    Thinker {-}{-}\textgreater{} Doer}
    
\NormalTok{    Doer {-}{-}\textgreater{} Merge\{\{Merge\}\}}
\NormalTok{    Merge {-}{-}\textgreater{} Coherence["Coherence Check (0.5)"]}
\NormalTok{    Coherence {-}{-}\textgreater{} End([End])}
    
\NormalTok{    style Knower fill:aliceblue,stroke:blue,stroke{-}width:2px}
\NormalTok{    style Thinker fill:lightgreen,stroke:green,stroke{-}width:2px}
\NormalTok{    style Doer fill:lightcoral,stroke:red,stroke{-}width:2px}
\NormalTok{    style Coherence fill:yellow,stroke:orange,stroke{-}width:2px}
\end{Highlighting}
\end{Shaded}

\textbf{Description:} Multi-scale reasoning operation flowchart showing
Knower Scale (2048D, exhaustive search), Thinker Scale (512D, trajectory
prediction), Doer Scale (12D, physical grounding), coherence check (0.5
target), and execution flow.

\begin{center}\rule{0.5\linewidth}{0.5pt}\end{center}

\subsection{Figure Index}\label{figure-index}

\begin{longtable}[]{@{}
  >{\raggedright\arraybackslash}p{(\linewidth - 6\tabcolsep) * \real{0.2162}}
  >{\raggedright\arraybackslash}p{(\linewidth - 6\tabcolsep) * \real{0.3514}}
  >{\raggedright\arraybackslash}p{(\linewidth - 6\tabcolsep) * \real{0.2162}}
  >{\raggedright\arraybackslash}p{(\linewidth - 6\tabcolsep) * \real{0.2162}}@{}}
\toprule\noalign{}
\begin{minipage}[b]{\linewidth}\raggedright
Figure
\end{minipage} & \begin{minipage}[b]{\linewidth}\raggedright
Description
\end{minipage} & \begin{minipage}[b]{\linewidth}\raggedright
Format
\end{minipage} & \begin{minipage}[b]{\linewidth}\raggedright
Source
\end{minipage} \\
\midrule\noalign{}
\endhead
\bottomrule\noalign{}
\endlastfoot
FIG. 1 & System Architecture & PNG/PDF/SVG &
\texttt{mermaid/fig01\_architecture.mmd} \\
FIG. 2 & VAE Architecture & PNG/PDF/SVG & \texttt{svg/fig02\_vae.svg} \\
FIG. 3 & 12D Physics State & PNG/PDF/SVG &
\texttt{svg/fig03\_12d\_state.svg} \\
FIG. 4 & Trajectory Prediction & PNG/PDF/SVG &
\texttt{svg/fig04\_trajectory.svg} \\
FIG. 5 & HIHO Double-Well & PNG/PDF & \texttt{python/plot\_hiho.py} \\
FIG. 6 & Training Convergence & PNG/PDF &
\texttt{python/plot\_training.py} \\
FIG. 7 & Journey Tracking & PNG/PDF/SVG &
\texttt{mermaid/fig07\_journey.mmd} \\
FIG. 8 & Multi-Scale Reasoning & PNG/PDF/SVG &
\texttt{mermaid/fig08\_multi\_scale.mmd} \\
\end{longtable}

\begin{center}\rule{0.5\linewidth}{0.5pt}\end{center}

\subsection{Reproducibility}\label{reproducibility}

All figures are reproducible from source files:

\begin{Shaded}
\begin{Highlighting}[]
\CommentTok{\# Render all figures}
\ExtensionTok{quarto}\NormalTok{ render figures.qmd }\AttributeTok{{-}{-}output}\NormalTok{ figures.pdf}

\CommentTok{\# Generate Mermaid PNGs}
\ExtensionTok{mmdc} \AttributeTok{{-}i}\NormalTok{ mermaid/fig01\_architecture.mmd }\AttributeTok{{-}o}\NormalTok{ png/fig01\_architecture.png }\AttributeTok{{-}w}\NormalTok{ 2000 }\AttributeTok{{-}H}\NormalTok{ 1200}
\ExtensionTok{mmdc} \AttributeTok{{-}i}\NormalTok{ mermaid/fig07\_journey.mmd }\AttributeTok{{-}o}\NormalTok{ png/fig07\_journey.png }\AttributeTok{{-}w}\NormalTok{ 2000 }\AttributeTok{{-}H}\NormalTok{ 1200}
\ExtensionTok{mmdc} \AttributeTok{{-}i}\NormalTok{ mermaid/fig08\_multi\_scale.mmd }\AttributeTok{{-}o}\NormalTok{ png/fig08\_multi\_scale.png }\AttributeTok{{-}w}\NormalTok{ 2000 }\AttributeTok{{-}H}\NormalTok{ 1200}

\CommentTok{\# Generate Python plots}
\ExtensionTok{python}\NormalTok{ python/plot\_hiho.py}
\ExtensionTok{python}\NormalTok{ python/plot\_training.py}
\end{Highlighting}
\end{Shaded}

\begin{center}\rule{0.5\linewidth}{0.5pt}\end{center}

\subsection{References}\label{references}

\begin{itemize}
\tightlist
\item
  Smith, W. B. (1962). ``The New Science'' (unpublished manuscript)
\item
  Percival, H. (1946). ``Thinking and Destiny''
\item
  Kingma, D. P., \& Welling, M. (2013). ``Auto-Encoding Variational
  Bayes''
\item
  Cohezion source code: github.com/manderson240/cohezion
\end{itemize}




\end{document}
